\chapter{Mạnh và Yếu}
\markboth{Mạnh và Yếu}{Strong and Weak}

\section{Một người nói rất cứng nhưng không dám nhận lỗi.}
\begin{truyen}{Một người nói rất cứng nhưng không dám nhận lỗi.}{Hardness is not always strength.}
\chuhoa{A}{nh quen tỏ ra mạnh mẽ, quyết liệt, không nhượng bộ. Nhưng mỗi khi sai, anh lại tìm cách tránh né. Về sau anh mới hiểu có những bề ngoài rất cứng mà bên trong lại sợ hãi trước sự thật rất nhỏ.}
\textit{(He was used to appearing tough, decisive, and unwilling to yield. But whenever he was wrong, he avoided admitting it. Later he understood that some exteriors are hard while the inside is afraid of even a small truth.)}
\end{truyen}
\begin{baihoc}
Mạnh không nằm ở vẻ gân guốc, mà ở chỗ dám đối diện phần sai của mình.
\textit{(Strength is not in rough appearance, but in facing the part of yourself that is wrong.)}
\end{baihoc}
\begin{ghinhoanh}
\textbf{Short English to remember:}

Hardness is not always strength.
\end{ghinhoanh}
\begin{tuhocnhanh}
1. Em có đang nhầm sự cứng với sự mạnh không? \textit{(Are you confusing hardness with strength?)}

2. Sự thật nào em còn ngại nhìn vào? \textit{(What truth are you still afraid to face?)}
\end{tuhocnhanh}
\ngancach

\section{Một người dám khóc sau nhiều năm gồng mình.}
\begin{truyen}{Một người dám khóc sau nhiều năm gồng mình.}{Letting down the guard can be strength.}
\chuhoa{A}{nh từng nghĩ không được phép mềm xuống. Nhưng khi thôi gồng và cho mình được yếu đi một chút, anh lại nhẹ hơn, thật hơn và gần người khác hơn. Có những giọt nước mắt không làm con người yếu đi; nó làm con người bớt giả vờ mạnh.}
\textit{(He once thought he was never allowed to soften. But when he stopped forcing himself and allowed some vulnerability, he became lighter, truer, and closer to others. Some tears do not weaken a person; they make the person stop pretending to be strong.)}
\end{truyen}
\begin{baihoc}
Dám thừa nhận mình mệt đôi khi mạnh hơn cố tỏ ra không sao.
\textit{(Admitting you are tired is sometimes stronger than pretending you are fine.)}
\end{baihoc}
\begin{ghinhoanh}
\textbf{Short English to remember:}

Vulnerability can be strength.
\end{ghinhoanh}
\begin{tuhocnhanh}
1. Lâu rồi em có cho mình được mềm xuống thật chưa? \textit{(How long has it been since you allowed yourself to soften honestly?)}

2. Em đang gồng ở chỗ nào quá lâu? \textit{(Where have you been forcing yourself for too long?)}
\end{tuhocnhanh}
\ngancach

\section{Một người nhìn hiền nhưng rất vững khi cần bảo vệ điều đúng.}
\begin{truyen}{Một người nhìn hiền nhưng rất vững khi cần bảo vệ điều đúng.}{Gentleness can carry backbone.}
\chuhoa{C}{hị nói nhỏ, cư xử nhẹ, nên nhiều người tưởng chị yếu. Nhưng khi thấy điều sai, chị vẫn giữ lập trường rất rõ. Lúc đó người ta mới nhận ra dịu dàng và vững vàng có thể đi cùng nhau.}
\textit{(She spoke softly and acted gently, so many thought she was weak. Yet when facing wrongdoing, she held a very clear stand. People then realized that gentleness and firmness can live together.)}
\end{truyen}
\begin{baihoc}
Yếu không nằm ở giọng nhỏ; yếu là khi biết điều đúng mà không dám giữ.
\textit{(Weakness is not in a soft voice; weakness is knowing what is right and not daring to hold it.)}
\end{baihoc}
\begin{ghinhoanh}
\textbf{Short English to remember:}

Gentleness can carry backbone.
\end{ghinhoanh}
\begin{tuhocnhanh}
1. Em có đang đánh giá sức mạnh của người khác quá nông không? \textit{(Are you judging others' strength too superficially?)}

2. Em muốn giữ sự vững vàng của mình theo cách nào? \textit{(How do you want to keep your own firmness?)}
\end{tuhocnhanh}
\ngancach

\section{Một đứa trẻ yếu hơn bạn bè nhưng chịu tập đều mỗi ngày.}
\begin{truyen}{Một đứa trẻ yếu hơn bạn bè nhưng chịu tập đều mỗi ngày.}{Weakness can become a teacher.}
\chuhoa{L}{úc đầu em làm gì cũng chậm hơn các bạn. Nhưng vì biết mình yếu hơn, em chịu khó tập đều, không bỏ. Nhiều năm sau, chính sự thiếu lợi thế ban đầu lại thành lý do em có sức bền hơn.}
\textit{(At first he was slower than others at everything. But because he knew he was weaker, he practiced steadily and did not quit. Years later, that early disadvantage became the reason he had more endurance.)}
\end{truyen}
\begin{baihoc}
Yếu thế lúc đầu không luôn là bất lợi; có khi nó buộc ta xây phần bền mà người mạnh sẵn không cần xây.
\textit{(Early weakness is not always a disadvantage; sometimes it forces us to build endurance that the naturally strong never had to build.)}
\end{baihoc}
\begin{ghinhoanh}
\textbf{Short English to remember:}

Weakness can become a teacher.
\end{ghinhoanh}
\begin{tuhocnhanh}
1. Điểm yếu nào của em đang có thể dạy em điều lớn hơn? \textit{(What weakness in you might be teaching a larger lesson?)}

2. Em đang né nó hay đang học từ nó? \textit{(Are you avoiding it or learning from it?)}
\end{tuhocnhanh}
\ngancach

\section{Một người mạnh với thiên hạ nhưng rất yếu trước sự cô đơn.}
\begin{truyen}{Một người mạnh với thiên hạ nhưng rất yếu trước sự cô đơn.}{Strength in public is not strength in every room.}
\chuhoa{B}{ên ngoài anh giải quyết việc rất giỏi, ai cũng thấy anh bản lĩnh. Nhưng khi đêm xuống và không còn công việc để trốn vào, anh lại thấy mình mong manh. Có những sức mạnh chỉ dùng được ngoài đời, chưa chạm tới phần sâu nhất bên trong.}
\textit{(Outwardly he handled life well and looked capable to everyone. But at night, when work could no longer distract him, he felt fragile. Some strengths only work in public life and have not yet reached the deepest part within.)}
\end{truyen}
\begin{baihoc}
Hiểu chỗ yếu sâu nhất của mình là bước đầu để mạnh lên thật sự.
\textit{(Understanding your deepest weakness is the first step toward real strength.)}
\end{baihoc}
\begin{ghinhoanh}
\textbf{Short English to remember:}

Public strength is not whole strength.
\end{ghinhoanh}
\begin{tuhocnhanh}
1. Phần yếu nhất của em nằm ở đâu? \textit{(Where is your deepest weakness?)}

2. Em đang che nó hay đang chăm nó? \textit{(Are you hiding it or caring for it?)}
\end{tuhocnhanh}
\ngancach